\chapter[Sermon 3. Of the Beginning of This Book, and the Apostles Salutation: Wherein Are Declared the Mysteries Chiefly of Christ, Secondly of Our Whole Faith and Redemption.]{Of the Beginning of This Book, and the Apostles Salutation: Wherein Are Declared the Mysteries Chiefly of Christ, Secondly of Our Whole Faith and Redemption.}
\label{sermon:003}
\markright{Sermon 3. Of the Beginning of This Book, and the Apostles ...}

John to the seven congregations in Asia. Grace be with you and peace from him who is, and who was, and who is to come, and from the seven spirits which are sent forth before his throne. And from Jesus Christ, the faithful witness, the firstborn of the dead, and Lord over the kings of the earth. To him who loved us and washed us from our sins in his own blood, and made us kings and priests to God, to him be glory and dominion forevermore. Amen. Behold, he comes with clouds, and every eye shall see him, and they also who pierced him. And all kindreds of the earth shall wail over him. Even so. Amen. I am Alpha and Omega, the beginning and the ending, says the Lord Almighty, who is, and who was, and who is to come.

Another section of the first part of this book contains a beginning or preface, wherein is the Apostle’s salutation, whereby he first describes the whole mystery of Christ, and secondly of our faith and redemption. For so the Apostles were accustomed in the beginning of their writings to include a brief summary of salutation. This practice is evident everywhere in Paul’s Epistles. By this description, he gains the goodwill and attention of all people.

\index{John, Saint}\index{Arethas of Caesarea}\index{Primasius}\index{Rome}The opening greeting of the Apostles is nothing else but a blessing. Blessing is an old, customary practice, by which the patriarchs wished for their children all manner of good things, both of body and soul. This is described at length in Genesis. Also, the high priest had authority given to bless the people, as we read in the sixth of Numbers, especially commanding that God’s name be placed upon the people. Therefore, it is a superstition to say that God, from whom every good gift descends from above, blesses, that is, gives good things, but that ministers or men merely wish. And the Lord in the law promises that he will grant those things to the people which the high priests should wish for them. Therefore, neither words nor shaven crowns, but the truth and power of God give the gifts. We ought not therefore to doubt that God will grant to us also the apostolic blessing, so that, being reconciled and accepted by God, we might have peace. And first, Saint John repeats his name, lest we should doubt of the author, whom we see Christ to have used as scribe and interpreter to all congregations. But he does not repeat that he is that servant of God, or witness, or Apostle of Jesus Christ. That was sufficient to have heard at the first beginning. Therefore, he teaches modesty and humility, which have received great gifts. Afterward, he indicates to whom he writes, and to whom this book pertains, to the seven churches of Asia, the names of which he will mention shortly after. Arethas, bishop of Caesarea, by the seven churches, and by the seven numbers, says he signified the multitude of churches that are in all places. So also Primasius, bishop of Utica in Africa, explains the seven numbers. Therefore, this greeting, this book, and the whole doctrine of Jesus Christ, written by Saint John, pertains to the whole universal church of Christ throughout all the world, and in all times and ages. Whereupon it belongs to all of us as many as are of us in the church of Christ. For although the letters are addressed to the Romans and Galatians, it does not follow that they are not ours. And he writes expressly to the churches of Asia, not to the churches of Jerusalem or Judaea, that he might show that the kingdom of Christ has already come to the Gentiles. And as God from the beginning chose Israel, in whom he might set forth a perfect example of the church and commonwealth, so from the beginning of the New Testament, he chose those seven churches of Asia, which he might set forth to the whole Christian world. But if Rome had been placed in the first place among the churches, as Ephesus is, good God, how much would the Roman faction make of it, for the establishing of their supremacy?

And the manner of the Apostles’ salutation wishes grace. Grace is the favor of God, and the reconciliation meant, whereby God the Father for Christ’s sake is made one with us, our sins pardoned, and we adopted for his children. Thereof arises the peace and tranquility of mind and the desire of concord with all men.

And here he shows abundantly who grants the church his blessing, that is to say, grace, reconciliation, and peace—God, and God three in persons, the Father, the Son, and the Holy Ghost, one God in essence. But here he discerns the persons very well. From him that is, to wit, the Father; and from the seven spirits, that is, from the Holy Ghost; and from Jesus Christ—this is the diversity of persons. And the signification of the unity is, when after the properties of persons are declared, he adds, “I am Alpha and Omega.” \&c. And that the Holy Ghost is set here in the midst, it disorders not the mystery of the Trinity, but appears to be an argument that he is the spirit as well of the Father as of the Son, and that he proceeds from both. As it is also proved by the words of our Lord, the 14th, 15th, and 16th of John. Here is also described the whole wholesome mystery first of Christ, then of the catholic faith, and of our redemption, so that herein you may find the chiefest articles of the Apostles' Creed, and have here a most goodly description of Christ our Lord. Hereof all men shall judge how truly some say that this book, contrary to the common opinion of the Apostles, makes little mention of Christ and of faith.

The father, as source and origin, from whom the Son is generated, is first described: for it is he who is, who was, and who is to come. John took those words from Moses in the third and thirty-fourth chapters of Exodus, and from many testimonies of Isaiah. And he says nothing but that God the Father is an eternal essence, which exists by and of itself, and is and preserves life in all. And that this essence is such that it has been always without beginning. For this is what he joins to being, or existence, “was.” He adds, and he who shall come, which shall be, and shall remain even to the end, and to everlastingness without end. The Greeks derive [illegible] from running, for that concerning and running, he meddles with all matters; is everywhere present, bringing help to the godly, or restraining and punishing the wicked.

And the Holy Spirit, where he is but one, for the sevenfold—that is, all manner of grace and manifold gifts—is here called, as I may say, Septenary or of the seventh number. And from the seven spirits, John says that it is from that Spirit, which is endowed with the sevenfold grace. Those diverse gifts are, in a way, explained in the eleventh chapter [?of Revelation] and elsewhere in the scriptures. He is said to be in the sight of the throne, that is, before the throne of God, joined truly in governance with the Father and the Son. For the throne is often used to represent the kingdom. Therefore, the Holy Spirit is of the same glory, power, and majesty as God.

Now he has come to Christ, whom by his attributes he describes most fully. You know that Jesus is the proper name of Christ, which Matthew explains as “Saviour.” Christ is the surname of his office and dignity, as one might say “anointed,” that is, bishop and king.

First he calls Christ our Lord a faithful witness, as in chapters 49 and 50 of Isaiah. For he was sent by the Father into the world from heaven, an Apostle, to testify God’s will, what he would have done with humankind—namely, that he would save the world by his Son and through faith in him, who is obedient to God’s law. For he must do the will of his Father. This Christ is a faithful witness, meaning sure, constant, and true; of whose doctrine no man ought to doubt. No one has seen God at any time; the only-begotten Son, who is in the Father’s embrace, has revealed him. This one, therefore, is the overseer and universal guide of the church. Whoever dissents from him are to be shunned. “Hear him,” says the Father.

\index{Resurrection}2 He is the firstborn of the dead. For he died indeed for our sins; he rose again from the dead, and was made the firstborn of the dead, Lord and conqueror of death. In him we believe that we also shall rise again, and in what manner. As Paul writes in 1 Corinthians 15. And just as in the first part he shadows the humanity of Christ, wherein he taught also his deity, in that he was the faithful, true, and catholic bishop, and is yet at this day; so in the second, the articles of our belief concerning the death of Christ and his resurrection are confirmed. To these may also be added the article of the resurrection of the dead.

3 Christ is prince over the kings of the earth, a monarch truly, and Lord of all rulers: He has taken a name above all names, the Lord of angels and of all creatures, to whom all things are subject, as the Apostle expounds in Colossians and Philippians 2. And he does not abolish laws and magistrates which will be king of kings and Lord of lords. For if there were no kings, how could Christ be king of kings? The most sacred Emperors, Constantine, Theodosius, and Justinian, knew themselves to be clients of Christ: their kingdom was Christ's, and they were to be subjects. Christ acknowledges these as his, by whom he governs those he has redeemed with his blood. Those who proudly rule over the people, boast themselves to be lords of all things, and do not acknowledge Christ to be monarch over all, are stark mad. And here are comprehended such things as we confess in the articles of our faith, that Christ ascended into heaven and sits on the right hand of the Father: that is, that he has received the highest power over all things in heaven and earth, as in Ephesians 1 and Acts 2.

4 Christ has loved us with incomparable love. For he himself says: “Greater love has no one than that someone should lay down his life for his friends.” The Apostle amplifies this in the fifth chapter to the Romans. And it was exceeding great love that moved Christ to come down from heaven and be incarnate and to redeem us by his death. With a free love he loved us, provoked by no merit of ours. For as John in his canonical Epistle speaks the same of the Father, “In this is love, not that we have loved God, but that he has loved us and sent his son as an atoning sacrifice for our sins,” so it is to be understood of the Son, that he has and does bear us great good will, not moved thereto through our love, with which we have embraced him. And of the free love to mankind, he gave himself unto death and washed us from our sins. For it is straightway added, by his blood.

\index{Augustine, Saint}Where three things seem to be observed by us. First, that Christ washes, purges, purifies, or cleanses the faithful: and that most fully, not partly. He alluded to the washings of the law, which he also expounded. For David says: “Purge me with hyssop, and I shall be made clean; wash me, and I shall be whiter than snow.” The same phrase of speech repeats Augustine in the first chapter. Micah also says: “The Lord will return and will have mercy on us; he will trample underfoot our iniquities. And you shall throw into the depths of the sea all their sins.” And the Lord says: “I will pour clean waters upon you, and you shall be cleansed from all your uncleanness.” The Lord Christ accomplishing these things, washes us, purges, and cleanses thoroughly, as well from the fault as the pain. He cleanses us from our sins, not from one, but from all. This which thing is proved both by former testimonies, and again in the first and second Epistle of John. Lastly, the manner of purifying is set forth, by blood. For without the shedding of blood no remission was made. Therefore through the mediation of death and bloodshedding, full remission of all sins was obtained for the faithful. They who bring forth any other manner of forgiveness of sins, are injurious to the death and blood of the Son of God. And here we may plainly see set forth an article of the Apostolic creed: I believe the forgiveness of sins.

\index{Zechariah (Book of)}In the fifth place, the effect of our redemption and purification is shown. For Christ has accomplished that as many of us as believe in the Father through the Son of God should be kings and priests to God and to his Father. Arethas and the copy of the complete text read not “kings” but “kingdom,” which is not read amiss. For we are the kingdom of God, because God by his Spirit, not the flesh nor the world, ought to reign in us. And when we submit to the Spirit’s governance, we are the kingdom of God. Which thing Paul discusses at length in chapter six to the Romans. Moreover, we are made kings—that is, free—by Christ, so that we should not serve the devil, the flesh, and the world, according to that saying of Zechariah, being delivered from the hands of our enemies, that we might serve him without fear in holiness and righteousness before him all the days of our life. And Christ has consecrated priests with his Spirit and blood, that we should offer up to God spiritual sacrifices, ourselves pure, prayers and praises, and almsdeeds. For that these are spiritual offerings, Peter and Paul testify. And John took these things from Exodus: For we of the Gentiles who have believed have succeeded in the place of the people of Israel, rejected Christ through unbelief. And these things give light to that article of the Creed, “I believe in the holy catholic church, the communion of saints.” For we, as many of us as believe, are the fellowship of God’s people, sanctified through Christ to the service of God. Of whom are these things said?

\index{Papacy}In the sixth place, in the description of Christ, he shows that glory and rule are due to God alone through Christ, for the church forevermore. We give glory to God when we ascribe our salvation and all goodness to his goodness, not to our own strength and merits. We give him rule when we acknowledge him to be Lord and head in the church, working by himself, not by the saints in heaven, to whom he has granted power. Not by the Pope, whom he has not constituted Vicar on earth. The whole glory and rule is Christ’s.

Seventhly, in the description follows the coming of Christ for judgment, and the manner of his coming. Just as a cloud took him from the sight of the Apostles, even so shall he come in clouds to judge the living and the dead. The scripture testifies. And he adds that the eyes of all men shall see the judge, even of those who have displeased him. From this we gather two things: first, that the judgment shall be universal, wherein men rising shall see Christ with their own eyes. Another thing is that Christ shall come to judge in the same flesh, in which he was wounded and pierced, hung upon the cross, was buried, and rose again. This passage is taken from Zechariah, and is also cited in Saint John’s Gospel. And it behooves that his body be shown to the whole world full of prints and marks, so that from this may be judged the godly and also the ungodly: those who then believed in such a redeemer, these who then rejected and scorned such a one. Of these we understand that is added: And they shall wail, for indeed they have neglected their own salvation, which the wise man discourses at large. Moreover, lest any should doubt about those things that are spoken of the judgment, and of the lamentation of the wicked (as Saint Peter said, the contemners and mockers of the judgment would be), he adds a kind of confirmation, even so. Amen.

In them also is explained the article of the creed of Christ who will judge the living and the dead. He concludes this passage with these words: “I am Alpha and Omega,” with the phrase “that which follows (the beginning and end)” omitted in some copies. It seems that interpretation of “I am Alpha and Omega” crept in from the margin. It is a saying of Saint John the Apostle, “I am Alpha and Omega.” Heretics, such as Basilides and Valentine, were greatly delighted in letters. But against those lettered heretics, John speaks plainly by the mouth of Christ: “I am Alpha and Omega.” If anything ought to be ascribed to letters, I am all of this—the whole everlasting virtue, essence, and eternity. For the meaning is that God is the beginning and end, that is, eternal, unspeakable, best, and greatest. These things are repeated: “He that is, which was…” and so on, which were explained before. There is added, “almighty.” For hereby is declared the unity and majesty of God, of whom the Trinity was also explained before. Hereby also the authority of this book is confirmed, the author of which is shown to be that eternal and almighty God. To whom be glory.
